% Imporditud: tools/import_eksperimendid.py
% Allikas: Eesti füüsikaolümpiaadi eksperimendi ülesannete
%   kogu (Rael Kalda, 2023), ülesanne Ü180, lahendus L180.
\setAuthor{Jarl Patrick Paide}
\setRound{piirkonnavoor}
\setYear{2022}
\setNumber{G E2}
\setDifficulty{5}
\setTopic{vedelike mehaanika}
\setVahendid{vesi anumaga, foolium, joonlaud, käärid}
\setPunktid{14}
\setAllikas{Eksperimendi kogu Ü180, lahendus lk 142}
\setLahendusseis{toores_ilma_skeemita}

\prob{Foolium}
Leidke fooliumi pindtihedus (mass pindala kohta). Vee pindpinevustegur $\sigma$ = 0,072 N m$^{-1}$ ja tihedus $\rho$ = 1000 kg m$^{-3}$.
\textit{Vihje:} vedelikupinna kaks mõttelist poolt tõmbavad üksteist jõuga F = $\sigma\ell$, kus $\ell$ on mõttelise eraldusjoone pikkus.

\hint

\solu
Üldine lahendusmeetod: Teeme fooliumist teatud kujuga objekti ning asetame selle vee peale. Tehes sobiliku kujuga objekti ning vajalikud mõõtmised saame leida fooliumi pindtiheduse Üldiselt kehtib vees olevale objektile 3 jõudu: üleslükkejõud, pindpinevusjõud ja raskusjõud. Enamasti on vaja leida ka fooliumi paksus. Selleks tuleb võtta võimalikult suur tükk fooliumit ning voltida/lõigata (voltimiste arv peab olema piisavalt väike, et see ei mõjutaks paksust) seda nii, et võimalikult palju kihte oleks üksteise peal. Kui on suurusjärgus 100 kihti, on võimalik joonlauaga mõõta paksust (küll mitte väga täpselt). Võib olla võimalik ka valides objekti kuju nii, et pindpinevusjõud on palju suurem kui üleslükkejõud. Sellisel juhul ei ole ilmselt vaja leida fooliumi pakusust, mille mõõtmine antud vahenditega on suhteliselt ebatäpne.

Fooliumi pindtihedus sõltub fooliumi paksusest ning on võrdne fooliumi tiheduse ja selle pakuses korrutisega. Fooliumi tihedus on 2700 kg m$^{-3}$. Sõltuvalt fooliumi pakusest on pindtihedus enamasti 20 g m$^{-2}$ kuni 80 g m$^{-2}$ (hindamisel oleks hea tea-

da fooliumi täpset pindtihedust).

Näiduslahendus:

Teeme fooliumist täidetud risttahuka mille alustahk on mõõtmetega a $\times$ b ja kõr-

guse moodustab N fooliumi paksust d. Keha uppumise piirjuhul kehtib võrrand

2(a + b)$\sigma$ + $\rho$v(abN d)g = $\rho$f (abN d)g. Võrrandist saame avaldada fooliumi pindti-

heduse: 2(a + b)$\sigma$ pindtihedus = $\rho$f d = abN g + $\rho$vd.

\textit{Žürii hindamisskeem on kogumikus lk 142; siia ei ole seda ümber kirjutatud.}
\probend
